\section{From the influence functional to the Hamiltonian of mean force}
\label{sec:hmf}

Equation~\eqref{eq:reduced_effective_PI} gives an exact reduced equilibrium path
integral after elimination of the bath. Its difficulty is not conceptual - the
bath has been integrated out - but structural: the influence functional produces
a \emph{bilocal} self-interaction in imaginary time. This nonlocality is encoded
by the translation-invariant kernel $K(\tau-\tau')$ determined by the bath
spectral density. In this form the reduced equilibrium object is explicit as a
functional integral, but it is not manifestly the kernel of a $\tau$-ordered
propagator generated by a local Hamiltonian on the system Hilbert space, and
therefore does not directly yield a tractable operator expression for
$H_{\mathrm{MF}}(\beta)=-(1/\beta)\log\bar\rho_Q(\beta)$.

The standard way to localise a Gaussian bilocal interaction is a
Hubbard--Stratonovich (HS) transformation. Applied here, HS replaces the bilocal
term by a linear coupling to an auxiliary Gaussian field. This step is the
bridge back from the influence-functional path integral to an operator-level
imaginary-time propagation picture. In particular, it produces exactly the
quenched canonical density introduced in Sec.~\ref{sec:quenched}.

\subsection{Hubbard--Stratonovich localisation of the Euclidean influence kernel}

For any symmetric positive kernel $K$ (as in \eqref{eq:kernel_spectral}) the
Gaussian identity
\begin{equation}
\begin{split}
    &\exp\!\bigg[
        \frac{1}{2}\int_0^\beta d\tau\int_0^\beta d\tau'
    \\
    &\qquad\quad\times
        j(\tau)\,K(\tau-\tau')\,j(\tau')
    \bigg]
    \\
    &=
    \mathcal N_K
    \int \mathcal D\xi\;
    \exp\bigg[
        -\frac{1}{2}\int_0^\beta d\tau
    \\
    &\qquad\quad\times
        \int_0^\beta d\tau'\,
        \xi(\tau)\,K^{-1}(\tau-\tau')\,\xi(\tau')
    \\
    &\qquad\qquad\qquad\quad
        +\int_0^\beta d\tau\,\xi(\tau)\,j(\tau)
    \bigg]
\end{split}
    \label{eq:HS_identity}
\end{equation}
holds, where $\mathcal N_K$ is a normalisation independent of $j$, and $K^{-1}$
denotes the inverse kernel on $\beta$-periodic functions. (Equivalently, one may
interpret the right-hand side as an expectation value over a centred Gaussian
field $\xi$ with covariance $K$:
$\langle\xi(\tau)\xi(\tau')\rangle = K(\tau-\tau')$.)

We apply \eqref{eq:HS_identity} to the bilocal influence functional in
\eqref{eq:reduced_effective_PI} with source
\begin{equation}
    j(\tau) := f\!\big(q(\tau)\big).
\end{equation}
This yields the identity, inside the system path integral,
\begin{align}
    &\exp\!\bigg[
        \frac{1}{2}\int_0^\beta d\tau\int_0^\beta d\tau'\,
        f\!\big(q(\tau)\big)
    \nonumber\\
    &\qquad\qquad\times
        K(\tau-\tau')\,f\!\big(q(\tau')\big)
    \bigg]
    \nonumber\\
    &\quad=
    \mathcal N_K
    \int \mathcal D\xi\;
    \exp\bigg[
        -\frac{1}{2}\int_0^\beta d\tau\int_0^\beta d\tau'\,
        \xi(\tau)
    \nonumber\\
    &\qquad\qquad\times
        K^{-1}(\tau-\tau')\,\xi(\tau')
    \nonumber\\
    &\qquad\qquad
        +\int_0^\beta d\tau\,\xi(\tau)\,f\!\big(q(\tau)\big)
    \bigg].
    \label{eq:HS_applied}
\end{align}
Substituting \eqref{eq:HS_applied} into \eqref{eq:reduced_effective_PI} and
interchanging the order of integration gives
\begin{align}
    \bar\rho_Q(q_f,q_i;\beta)
    &=
    Z_X(\beta)\,\mathcal N_K
    \int \mathcal D\xi\;
    \exp\bigg[
        -\frac{1}{2}\int_0^\beta d\tau
    \nonumber\\
    &\qquad\qquad\times
        \int_0^\beta d\tau'\,
        \xi(\tau)\,K^{-1}(\tau-\tau')\,\xi(\tau')
    \bigg]
    \nonumber\\
    &\quad\times
    \int_{q(0)=q_i}^{q(\beta)=q_f}\!\!\mathcal Dq\;
    \exp\bigg[
        -S_Q[q] 
    \nonumber\\
    &\qquad\qquad
        + \int_0^\beta d\tau\,\xi(\tau)\,f\!\big(q(\tau)\big)
    \bigg].
    \label{eq:reduced_kernel_after_HS}
\end{align}
The inner path integral is now local in imaginary time for fixed $\xi(\tau)$.

\subsection{Return to the quenched canonical density}

The inner functional integral in \eqref{eq:reduced_kernel_after_HS} is precisely
the Euclidean path-integral kernel of a quenched canonical density generated by
a $\tau$-dependent Hamiltonian. Comparing \eqref{eq:reduced_kernel_after_HS} with
the general construction in Sec.~\ref{sec:quenched}, we identify the quenched
generator
\begin{equation}
    \bar H(\tau;\xi)=H_Q-\xi(\tau)\,f,
    \label{eq:quenched_generator_identification}
\end{equation}
and therefore the corresponding quenched operator
\begin{equation}
    \tilde\rho_Q(\beta;\xi)
    :=
    \hat\tau\exp\!\left[-\int_0^\beta d\tau\,\big(H_Q-\xi(\tau)f\big)\right],
    \label{eq:tilde_rho_def_again}
\end{equation}
with kernel given by the local (for fixed $\xi$) path integral derived previously.
Equation~\eqref{eq:reduced_kernel_after_HS} may then be rewritten compactly as
\begin{gather}
    \bar\rho_Q(\beta)
    =
    Z_X(\beta)\,\mathbb E_\xi\!\left[\tilde\rho_Q(\beta;\xi)\right],
    \label{eq:bar_rho_as_quenched_average_final}
    \\
    \mathbb E_\xi[\xi(\tau)\xi(\tau')]=K(\tau-\tau'), \notag
\end{gather}
where $\mathbb E_\xi[\cdot]$ denotes the Gaussian functional average with
covariance $K$. This is the promised bridge: the nonlocal influence-functional
path integral is exactly equivalent to an average over quenched, system-local
imaginary-time propagators.

\subsection{Exact evaluation in the commuting case}

We now specialise to the commuting (QND) sector,
\begin{equation}
    [H_Q,f]=0.
    \label{eq:commuting_here}
\end{equation}
In this case $H_Q$ and $f$ commute, so the generator \eqref{eq:quenched_generator_identification}
commutes with itself at different imaginary times. The $\hat\tau$ ordering becomes
redundant and \eqref{eq:tilde_rho_def_again} factorises:
\begin{align}
    \tilde\rho_Q(\beta;\xi)
    &=
    \exp\!\left[-\int_0^\beta d\tau\,H_Q\right]\,
    \exp\!\left[+\int_0^\beta d\tau\,\xi(\tau)\,f\right]
    \nonumber\\
    &=
    e^{-\beta H_Q}\,
    \exp\!\left[f\,X\right],
    \qquad
    X:=\int_0^\beta d\tau\,\xi(\tau).
    \label{eq:tilde_rho_commuting_factorised}
\end{align}
Since $\xi$ is Gaussian and $X$ is a linear functional of $\xi$, the random variable
$X$ is Gaussian with mean zero and variance
\begin{align}
    \mathbb E_\xi[X]&=0, \nonumber\\
    \mathbb E_\xi[X^2]
    &= \int_0^\beta d\tau\int_0^\beta d\tau'\,K(\tau-\tau')
    \nonumber\\
    &\equiv C(\beta).
    \label{eq:Cbeta_def}
\end{align}
Therefore the Gaussian average can be performed exactly:
\begin{equation}
    \mathbb E_\xi\!\left[e^{f X}\right]
    =
    \exp\!\left[\frac{1}{2}C(\beta)\,f^2\right],
    \label{eq:gaussian_average_resummed}
\end{equation}
and hence, using \eqref{eq:bar_rho_as_quenched_average_final},
\begin{equation}
    \bar\rho_Q(\beta)
    =
    Z_X(\beta)\,e^{-\beta H_Q}\,\exp\!\left[\frac{1}{2}C(\beta)\,f^2\right].
    \label{eq:bar_rho_commuting_closed}
\end{equation}
Up to a $\beta$-dependent scalar fixed by the chosen normalisation, the Hamiltonian
of mean force is therefore
\begin{equation}
    H_{\mathrm{MF}}(\beta)
    =
    H_Q
    -\frac{C(\beta)}{2\beta}\,f^2
    +\mathrm{const}(\beta)\,\mathbf 1.
    \label{eq:HMF_commuting_final}
\end{equation}
The function $C(\beta)$ is determined entirely by the bath kernel, equivalently by
the spectral density via \eqref{eq:kernel_spectral}:
\begin{equation}
    C(\beta)=\int_0^\beta d\tau\int_0^\beta d\tau'\,K(\tau-\tau').
    \label{eq:Cbeta_kernel_moments}
\end{equation}
This completes the exact construction of $H_{\mathrm{MF}}(\beta)$ in the commuting sector.
